% Options for packages loaded elsewhere
\PassOptionsToPackage{unicode}{hyperref}
\PassOptionsToPackage{hyphens}{url}
\documentclass[
]{article}
\usepackage{xcolor}
\usepackage{amsmath,amssymb}
\setcounter{secnumdepth}{-\maxdimen} % remove section numbering
\usepackage{iftex}
\ifPDFTeX
  \usepackage[T1]{fontenc}
  \usepackage[utf8]{inputenc}
  \usepackage{textcomp} % provide euro and other symbols
\else % if luatex or xetex
  \usepackage{unicode-math} % this also loads fontspec
  \defaultfontfeatures{Scale=MatchLowercase}
  \defaultfontfeatures[\rmfamily]{Ligatures=TeX,Scale=1}
\fi
\usepackage{lmodern}
\ifPDFTeX\else
  % xetex/luatex font selection
\fi
% Use upquote if available, for straight quotes in verbatim environments
\IfFileExists{upquote.sty}{\usepackage{upquote}}{}
\IfFileExists{microtype.sty}{% use microtype if available
  \usepackage[]{microtype}
  \UseMicrotypeSet[protrusion]{basicmath} % disable protrusion for tt fonts
}{}
\makeatletter
\@ifundefined{KOMAClassName}{% if non-KOMA class
  \IfFileExists{parskip.sty}{%
    \usepackage{parskip}
  }{% else
    \setlength{\parindent}{0pt}
    \setlength{\parskip}{6pt plus 2pt minus 1pt}}
}{% if KOMA class
  \KOMAoptions{parskip=half}}
\makeatother
\usepackage{color}
\usepackage{fancyvrb}
\newcommand{\VerbBar}{|}
\newcommand{\VERB}{\Verb[commandchars=\\\{\}]}
\DefineVerbatimEnvironment{Highlighting}{Verbatim}{commandchars=\\\{\}}
% Add ',fontsize=\small' for more characters per line
\newenvironment{Shaded}{}{}
\newcommand{\AlertTok}[1]{\textcolor[rgb]{1.00,0.00,0.00}{\textbf{#1}}}
\newcommand{\AnnotationTok}[1]{\textcolor[rgb]{0.38,0.63,0.69}{\textbf{\textit{#1}}}}
\newcommand{\AttributeTok}[1]{\textcolor[rgb]{0.49,0.56,0.16}{#1}}
\newcommand{\BaseNTok}[1]{\textcolor[rgb]{0.25,0.63,0.44}{#1}}
\newcommand{\BuiltInTok}[1]{\textcolor[rgb]{0.00,0.50,0.00}{#1}}
\newcommand{\CharTok}[1]{\textcolor[rgb]{0.25,0.44,0.63}{#1}}
\newcommand{\CommentTok}[1]{\textcolor[rgb]{0.38,0.63,0.69}{\textit{#1}}}
\newcommand{\CommentVarTok}[1]{\textcolor[rgb]{0.38,0.63,0.69}{\textbf{\textit{#1}}}}
\newcommand{\ConstantTok}[1]{\textcolor[rgb]{0.53,0.00,0.00}{#1}}
\newcommand{\ControlFlowTok}[1]{\textcolor[rgb]{0.00,0.44,0.13}{\textbf{#1}}}
\newcommand{\DataTypeTok}[1]{\textcolor[rgb]{0.56,0.13,0.00}{#1}}
\newcommand{\DecValTok}[1]{\textcolor[rgb]{0.25,0.63,0.44}{#1}}
\newcommand{\DocumentationTok}[1]{\textcolor[rgb]{0.73,0.13,0.13}{\textit{#1}}}
\newcommand{\ErrorTok}[1]{\textcolor[rgb]{1.00,0.00,0.00}{\textbf{#1}}}
\newcommand{\ExtensionTok}[1]{#1}
\newcommand{\FloatTok}[1]{\textcolor[rgb]{0.25,0.63,0.44}{#1}}
\newcommand{\FunctionTok}[1]{\textcolor[rgb]{0.02,0.16,0.49}{#1}}
\newcommand{\ImportTok}[1]{\textcolor[rgb]{0.00,0.50,0.00}{\textbf{#1}}}
\newcommand{\InformationTok}[1]{\textcolor[rgb]{0.38,0.63,0.69}{\textbf{\textit{#1}}}}
\newcommand{\KeywordTok}[1]{\textcolor[rgb]{0.00,0.44,0.13}{\textbf{#1}}}
\newcommand{\NormalTok}[1]{#1}
\newcommand{\OperatorTok}[1]{\textcolor[rgb]{0.40,0.40,0.40}{#1}}
\newcommand{\OtherTok}[1]{\textcolor[rgb]{0.00,0.44,0.13}{#1}}
\newcommand{\PreprocessorTok}[1]{\textcolor[rgb]{0.74,0.48,0.00}{#1}}
\newcommand{\RegionMarkerTok}[1]{#1}
\newcommand{\SpecialCharTok}[1]{\textcolor[rgb]{0.25,0.44,0.63}{#1}}
\newcommand{\SpecialStringTok}[1]{\textcolor[rgb]{0.73,0.40,0.53}{#1}}
\newcommand{\StringTok}[1]{\textcolor[rgb]{0.25,0.44,0.63}{#1}}
\newcommand{\VariableTok}[1]{\textcolor[rgb]{0.10,0.09,0.49}{#1}}
\newcommand{\VerbatimStringTok}[1]{\textcolor[rgb]{0.25,0.44,0.63}{#1}}
\newcommand{\WarningTok}[1]{\textcolor[rgb]{0.38,0.63,0.69}{\textbf{\textit{#1}}}}
\setlength{\emergencystretch}{3em} % prevent overfull lines
\providecommand{\tightlist}{%
  \setlength{\itemsep}{0pt}\setlength{\parskip}{0pt}}
\usepackage{bookmark}
\IfFileExists{xurl.sty}{\usepackage{xurl}}{} % add URL line breaks if available
\urlstyle{same}
\hypersetup{
  hidelinks,
  pdfcreator={LaTeX via pandoc}}

\author{}
\date{}

\begin{document}

\section{Practical Byzantine Fault
Tolerance}\label{practical-byzantine-fault-tolerance}

Miguel Castro and Barbara Liskov Laboratory for Computer Science,
Massachusetts Institute of Technology, 545 Technology Square, Cambridge,
MA 02139 \{castro, liskov\}@lcs.mit.edu

\section{Abstract}\label{abstract}

This paper describes a new replication algorithm that is able to
tolerate Byzantine faults. We believe that Byzantine-fault-tolerant
algorithms will be increasingly important in the future because
malicious attacks and software errors are increasingly common and can
cause faulty nodes to exhibit arbitrary behavior. Whereas previous
algorithms assumed a synchronous system or were too slow to be used in
practice, the algorithm described in this paper is practical: it works
in asynchronous environments like the Internet and incorporates several
important optimizations that improve the response time of previous
algorithms by more than an order of magnitude. We implemented a
Byzantine-fault-tolerant NFS service using our algorithm and measured
its performance. The results show that our service is only 3\% slower
than a standard unreplicated NFS.

\section{1 Introduction}\label{introduction}

Malicious attacks and software errors are increasingly common. The
growing reliance of industry and government on online information
services makes malicious attacks more attractive and makes the
consequences of successful attacks more serious. In addition, the number
of software errors is increasing due to the growth in size and
complexity of software. Since malicious attacks and software errors can
cause faulty nodes to exhibit Byzantine (i.e., arbitrary) behavior,
Byzantine-fault-tolerant algorithms are increasingly important.

This paper presents a new, practical algorithm for state machine
replication {[}17, 34{]} that tolerates Byzantine faults. The algorithm
offers both liveness and safety provided at most
\(\left\lfloor\frac{n-1}{3}\right\rfloor\) out of a total of n replicas
are simultaneously faulty. This means that clients eventually receive
replies to their requests and those replies are correct according to
linearizability {[}14, 4{]}. The algorithm works in asynchronous systems
like the Internet and it incorporates important optimizations that
enable it to perform efficiently.

There is a significant body of work on agreement and replication
techniques that tolerate Byzantine faults (starting with {[}19{]}).
However, most earlier work (e.g., {[}3, 24, 10{]}) either concerns
techniques designed to demonstrate theoretical feasibility that are too
inefficient to be used in practice, or assumes synchrony, i.e., relies
on known bounds on message delays and process speeds. The systems
closest to ours, Rampart {[}30{]} and SecureRing {[}16{]}, were designed
to be practical, but they rely on the synchrony assumption for
correctness, which is dangerous in the presence of malicious attacks. An
attacker may compromise the safety of a service by delaying non-faulty
nodes or the communication between them until they are tagged as faulty
and excluded from the replica group. Such a denial-of-service attack is
generally easier than gaining control over a non-faulty node.

Our algorithm is not vulnerable to this type of attack because it does
not rely on synchrony for safety. In addition, it improves the
performance of Rampart and SecureRing by more than an order of magnitude
as explained in Section 7. It uses only one message round trip to
execute read-only operations and two to execute read-write operations.
Also, it uses an efficient authentication scheme based on message
authentication codes during normal operation; public-key cryptography,
which was cited as the major latency {[}29{]} and throughput {[}22{]}
bottleneck in Rampart, is used only when there are faults.

To evaluate our approach, we implemented a replication library and used
it to implement a real service: a Byzantine-fault-tolerant distributed
file system that supports the NFS protocol. We used the Andrew benchmark
\([15]\) to evaluate the performance of our system. The results show
that our system is only 3\% slower than the standard NFS daemon in the
Digital Unix kernel during normal-case operation.

Thus, the paper makes the following contributions:

\begin{itemize}
\item
  It describes the first state-machine replication protocol that
  correctly survives Byzantine faults in asynchronous networks.\\
\item
  It describes a number of important optimizations that allow the
  algorithm to perform well so that it can be used in real systems.
\item
  It describes the implementation of a Byzantine-fault-tolerant
  distributed file system.\\
\item
  It provides experimental results that quantify the cost of the
  replication technique.
\end{itemize}

The remainder of the paper is organized as follows. We begin by
describing our system model, including our failure assumptions. Section
3 describes the problem solved by the algorithm and states correctness
conditions. The algorithm is described in Section 4 and some important
optimizations are described in Section 5. Section 6 describes our
replication library and how we used it to implement a
Byzantine-fault-tolerant NFS. Section 7 presents the results of our
experiments. Section 8 discusses related work. We conclude with a
summary of what we have accomplished and a discussion of future research
directions.

\section{2 System Model}\label{system-model}

We assume an asynchronous distributed system where nodes are connected
by a network. The network may fail to deliver messages, delay them,
duplicate them, or deliver them out of order.

We use a Byzantine failure model, i.e., faulty nodes may behave
arbitrarily, subject only to the restriction mentioned below. We assume
independent node failures. For this assumption to be true in the
presence of malicious attacks, some steps need to be taken, e.g., each
node should run different implementations of the service code and
operating system and should have a different root password and a
different administrator. It is possible to obtain different
implementations from the same code base {[}28{]} and for low degrees of
replication one can buy operating systems from different vendors.
N-version programming, i.e., different teams of programmers produce
different implementations, is another option for some services.

We use cryptographic techniques to prevent spoofing and replays and to
detect corrupted messages. Our messages contain public-key signatures
{[}33{]}, message authentication codes {[}36{]}, and message digests
produced by collision-resistant hash functions {[}32{]}. We denote a
message \(m\) signed by node \(i\) as \(\langle m \rangle_{\sigma_i}\)
and the digest of message \(m\) by \(D(m)\) . We follow the common
practice of signing a digest of a message and appending it to the
plaintext of the message rather than signing the full message
\((\langle m \rangle_{\sigma_i}\) should be interpreted in this way).
All replicas know the others' public keys to verify signatures.

We allow for a very strong adversary that can coordinate faulty nodes,
delay communication, or delay correct nodes in order to cause the most
damage to the replicated service. We do assume that the adversary cannot
delay correct nodes indefinitely. We also assume that the adversary (and
the faulty nodes it controls) are computationally bound so that (with
very high probability) it is unable to subvert the cryptographic
techniques mentioned above. For example, the adversary cannot produce a
valid signature of a non-faulty node, compute the information summarized
by a digest from the digest, or find two messages with the same digest.
The cryptographic techniques we use are thought to have these properties
{[}33, 36, 32{]}.

\section{3 Service Properties}\label{service-properties}

Our algorithm can be used to implement any deterministic replicated
service with a state and some operations. The operations are not
restricted to simple reads or writes of portions of the service state;
they can perform arbitrary deterministic computations using the state
and operation arguments. Clients issue requests to the replicated
service to invoke operations and block waiting for a reply. The
replicated service is implemented by n replicas. Clients and replicas
are non-faulty if they follow the algorithm in Section 4 and if no
attacker can forge their signature.

The algorithm provides both safety and liveness assuming no more than
\(\left\lfloor\frac{n-1}{3}\right\rfloor\) replicas are faulty. Safety
means that the replicated service satisfies linearizability {[}14{]}
(modified to account for Byzantine-faulty clients {[}4{]}): it behaves
like a centralized implementation that executes operations atomically
one at a time. Safety requires the bound on the number of faulty
replicas because a faulty replica can behave arbitrarily, e.g., it can
destroy its state.

Safety is provided regardless of how many faulty clients are using the
service (even if they collude with faulty replicas): all operations
performed by faulty clients are observed in a consistent way by
non-faulty clients. In particular, if the service operations are
designed to preserve some invariants on the service state, faulty
clients cannot break those invariants.

The safety property is insufficient to guard against faulty clients,
e.g., in a file system a faulty client can write garbage data to some
shared file. However, we limit the amount of damage a faulty client can
do by providing access control: we authenticate clients and deny access
if the client issuing a request does not have the right to invoke the
operation. Also, services may provide operations to change the access
permissions for a client. Since the algorithm ensures that the effects
of access revocation operations are observed consistently by all
clients, this provides a powerful mechanism to recover from attacks by
faulty clients.

The algorithm does not rely on synchrony to provide safety. Therefore,
it must rely on synchrony to provide liveness; otherwise it could be
used to implement consensus in an asynchronous system, which is not
possible {[}9{]}. We guarantee liveness, i.e., clients eventually
receive replies to their requests, provided at most
\(\left\lfloor\frac{n-1}{3}\right\rfloor\) replicas are faulty and
\(delay(t)\) does not

grow faster than t indefinitely. Here, \(\text{delay}(t)\) is the time
between the moment t when a message is sent for the first time and the
moment when it is received by its destination (assuming the sender keeps
retransmitting the message until it is received). (A more precise
definition can be found in {[}4{]}.) This is a rather weak synchrony
assumption that is likely to be true in any real system provided network
faults are eventually repaired, yet it enables us to circumvent the
impossibility result in {[}9{]}.

The resiliency of our algorithm is optimal: \(3f+1\) is the minimum
number of replicas that allow an asynchronous system to provide the
safety and liveness properties when up to f replicas are faulty (see
{[}2{]} for a proof). This many replicas are needed because it must be
possible to proceed after communicating with n - f replicas, since f
replicas might be faulty and not responding. However, it is possible
that the f replicas that did not respond are not faulty and, therefore,
f of those that responded might be faulty. Even so, there must still be
enough responses that those from non-faulty replicas outnumber those
from faulty ones, i.e., n - 2f \textgreater{} f.~Therefore n
\textgreater{} 3f.

The algorithm does not address the problem of fault-tolerant privacy: a
faulty replica may leak information to an attacker. It is not feasible
to offer fault-tolerant privacy in the general case because service
operations may perform arbitrary computations using their arguments and
the service state; replicas need this information in the clear to
execute such operations efficiently. It is possible to use secret
sharing schemes {[}35{]} to obtain privacy even in the presence of a
threshold of malicious replicas {[}13{]} for the arguments and portions
of the state that are opaque to the service operations. We plan to
investigate these techniques in the future.

\section{4 The Algorithm}\label{the-algorithm}

Our algorithm is a form of state machine replication {[}17, 34{]}: the
service is modeled as a state machine that is replicated across
different nodes in a distributed system. Each state machine replica
maintains the service state and implements the service operations. We
denote the set of replicas by R and identify each replica using an
integer in \(\{0,\ldots,|\mathcal{R}|-1\}\) . For simplicity, we assume
\(|\mathcal{R}|=3f+1\) where f is the maximum number of replicas that
may be faulty; although there could be more than \(3f+1\) replicas, the
additional replicas degrade performance (since more and bigger messages
are being exchanged) without providing improved resiliency.

The replicas move through a succession of configurations called views.
In a view one replica is the primary and the others are backups. Views
are numbered consecutively. The primary of a view is replica p such that
\(p = v \mod |\mathcal{R}|\) , where v is the view number. View changes
are carried out when it appears that the primary has failed. Viewstamped
Replication {[}26{]} and Paxos {[}18{]} used a similar approach to
tolerate benign faults (as discussed in Section 8.)

The algorithm works roughly as follows:

\begin{enumerate}
\def\labelenumi{\arabic{enumi}.}
\tightlist
\item
  A client sends a request to invoke a service operation to the
  primary\\
\item
  The primary multicasts the request to the backups\\
\item
  Replicas execute the request and send a reply to the client\\
\item
  The client waits for \(f + 1\) replies from different replicas with
  the same result; this is the result of the operation.
\end{enumerate}

Like all state machine replication techniques {[}34{]}, we impose two
requirements on replicas: they must be deterministic (i.e., the
execution of an operation in a given state and with a given set of
arguments must always produce the same result) and they must start in
the same state. Given these two requirements, the algorithm ensures the
safety property by guaranteeing that all non-faulty replicas agree on a
total order for the execution of requests despite failures.

The remainder of this section describes a simplified version of the
algorithm. We omit discussion of how nodes recover from faults due to
lack of space. We also omit details related to message retransmissions.
Furthermore, we assume that message authentication is achieved using
digital signatures rather than the more efficient scheme based on
message authentication codes; Section 5 discusses this issue further. A
detailed formalization of the algorithm using the I/O automaton model
{[}21{]} is presented in {[}4{]}.

\section{4.1 The Client}\label{the-client}

A client c requests the execution of state machine operation o by
sending a \(\langle REQUEST, o, t, c \rangle_{\sigma_{c}}\) message to
the primary. Timestamp t is used to ensure exactly-once semantics for
the execution of client requests. Timestamps for c's requests are
totally ordered such that later requests have higher timestamps than
earlier ones; for example, the timestamp could be the value of the
client's local clock when the request is issued.

Each message sent by the replicas to the client includes the current
view number, allowing the client to track the view and hence the current
primary. A client sends a request to what it believes is the current
primary using a point-to-point message. The primary atomically
multicasts the request to all the backups using the protocol described
in the next section.

A replica sends the reply to the request directly to the client. The
reply has the form \(\langle\text{REPLY},v,t,c,i,r\rangle_{\sigma_{i}}\)
where v is the current view number, t is the timestamp of the
corresponding request, i is the replica number, and r is the result of
executing the requested operation.

The client waits for \(f + 1\) replies with valid signatures from
different replicas, and with the same t and r, before

accepting the result r. This ensures that the result is valid, since at
most f replicas can be faulty.

If the client does not receive replies soon enough, it broadcasts the
request to all replicas. If the request has already been processed, the
replicas simply re-send the reply; replicas remember the last reply
message they sent to each client. Otherwise, if the replica is not the
primary, it relays the request to the primary. If the primary does not
multicast the request to the group, it will eventually be suspected to
be faulty by enough replicas to cause a view change.

In this paper we assume that the client waits for one request to
complete before sending the next one. But we can allow a client to make
asynchronous requests, yet preserve ordering constraints on them.

\section{4.2 Normal-Case Operation}\label{normal-case-operation}

The state of each replica includes the state of the service, a message
log containing messages the replica has accepted, and an integer
denoting the replica's current view. We describe how to truncate the log
in Section 4.3.

When the primary, p, receives a client request, m, it starts a
three-phase protocol to atomically multicast the request to the
replicas. The primary starts the protocol immediately unless the number
of messages for which the protocol is in progress exceeds a given
maximum. In this case, it buffers the request. Buffered requests are
multicast later as a group to cut down on message traffic and CPU
overheads under heavy load; this optimization is similar to a group
commit in transactional systems {[}11{]}. For simplicity, we ignore this
optimization in the description below.

The three phases are pre-prepare, prepare, and commit. The pre-prepare
and prepare phases are used to totally order requests sent in the same
view even when the primary, which proposes the ordering of requests, is
faulty. The prepare and commit phases are used to ensure that requests
that commit are totally ordered across views.

In the pre-prepare phase, the primary assigns a sequence number, n, to
the request, multicasts a pre-prepare message with m piggybacked to all
the backups, and appends the message to its log. The message has the
form
\(\langle\langle\text{PRE-PREPARE},v,n,d\rangle_{\sigma_{p}},m\rangle\)
, where v indicates the view in which the message is being sent, m is
the client's request message, and d is m's digest.

Requests are not included in pre-prepare messages to keep them small.
This is important because pre-prepare messages are used as a proof that
the request was assigned sequence number n in view v in view changes.
Additionally, it decouples the protocol to totally order requests from
the protocol to transmit the request to the replicas; allowing us to use
a transport optimized for small messages for protocol messages and a
transport optimized for large messages for large requests.

A backup accepts a pre-prepare message provided:

\begin{itemize}
\tightlist
\item
  the signatures in the request and the pre-prepare message are correct
  and \(d\) is the digest for \(m\) ;\\
\item
  it is in view \(v\) ;\\
\item
  it has not accepted a pre-prepare message for view \(v\) and sequence
  number \(n\) containing a different digest;\\
\item
  the sequence number in the pre-prepare message is between a low water
  mark, h, and a high water mark, H.
\end{itemize}

The last condition prevents a faulty primary from exhausting the space
of sequence numbers by selecting a very large one. We discuss how H and
h advance in Section 4.3.

If backup i accepts the
\(\langle\langle\text{PRE-PREPARE},v,n,d\rangle_{\sigma_{p}},m\rangle\)
message, it enters the prepare phase by multicasting a
\(\langle\text{PREPARE},v,n,d,i\rangle_{\sigma_{i}}\) message to all
other replicas and adds both messages to its log. Otherwise, it does
nothing.

A replica (including the primary) accepts prepare messages and adds them
to its log provided their signatures are correct, their view number
equals the replica's current view, and their sequence number is between
\(h\) and \(H\) .

We define the predicate \(\text{prepared}(m, v, n, i)\) to be true if
and only if replica i has inserted in its log: the request m, a
pre-prepare for m in view v with sequence number n, and 2f prepares from
different backups that match the pre-prepare. The replicas verify
whether the prepares match the pre-prepare by checking that they have
the same view, sequence number, and digest.

The pre-prepare and prepare phases of the algorithm guarantee that
non-faulty replicas agree on a total order for the requests within a
view. More precisely, they ensure the following invariant: if
\(\text{prepared}(m, v, n, i)\) is true then
\(\text{prepared}(m', v, n, j)\) is false for any non-faulty replica j
(including i = j) and any \(m'\) such that \(D(m') \neq D(m)\) . This is
true because \(\text{prepared}(m, v, n, i)\) and
\(|\mathcal{R}| = 3f + 1\) imply that at least \(f + 1\) non-faulty
replicas have sent a pre-prepare or prepare for m in view v with
sequence number n.~Thus, for \(\text{prepared}(m', v, n, j)\) to be true
at least one of these replicas needs to have sent two conflicting
prepares (or pre-prepares if it is the primary for v), i.e., two
prepares with the same view and sequence number and a different digest.
But this is not possible because the replica is not faulty. Finally, our
assumption about the strength of message digests ensures that the
probability that \(m \neq m'\) and \(D(m) = D(m')\) is negligible.

Replica \(i\) multicasts a
\(\langle \text{COMMIT}, v, n, D(m), i \rangle_{\sigma_i}\) to the other
replicas when prepared \((m, v, n, i)\) becomes true. This starts the
commit phase. Replicas accept commit messages and insert them in their
log provided they are properly signed, the view number in the message is
equal to the replica's current view, and the sequence number is between
\(h\) and \(H\)

We define the committed and committed-local predicates as follows:
committed(m,v,n) is true if and only if prepared(m,v,n,i) is true for
all i in some set of f+1 non-faulty replicas; and
committed-local(m,v,n,i) is true if and only if prepared(m,v,n,i) is
true and i has accepted \(2f+1\) commits (possibly including its own)
from different replicas that match the pre-prepare for m; a commit
matches a pre-prepare if they have the same view, sequence number, and
digest.

The commit phase ensures the following invariant: if
committed-local(m,v,n,i) is true for some non-faulty i then
committed(m,v,n) is true. This invariant and the view-change protocol
described in Section 4.4 ensure that non-faulty replicas agree on the
sequence numbers of requests that commit locally even if they commit in
different views at each replica. Furthermore, it ensures that any
request that commits locally at a non-faulty replica will commit at
\(f + 1\) or more non-faulty replicas eventually.

Each replica i executes the operation requested by m after
committed-local(m, v, n, i) is true and i's state reflects the
sequential execution of all requests with lower sequence numbers. This
ensures that all non-faulty replicas execute requests in the same order
as required to provide the safety property. After executing the
requested operation, replicas send a reply to the client. Replicas
discard requests whose timestamp is lower than the timestamp in the last
reply they sent to the client to guarantee exactly-once semantics.

We do not rely on ordered message delivery, and therefore it is possible
for a replica to commit requests out of order. This does not matter
since it keeps the preprepare, prepare, and commit messages logged until
the corresponding request can be executed.

Figure 1 shows the operation of the algorithm in the normal case of no
primary faults. Replica 0 is the primary, replica 3 is faulty, and C is
the client.

\textit{[Image omitted: images/2f4f7739a1d51a88e8375de4bd0e94768eb824f0a4fd50c39be7c98ab5bd543a.jpg]}

flowchart

Request processing flowchart with five stages: request, pre-prepare,
prepare, commit, reply, showing bidirectional and unidirectional message
exchanges.

Figure 1: Normal Case Operation

\section{4.3 Garbage Collection}\label{garbage-collection}

This section discusses the mechanism used to discard messages from the
log. For the safety condition to hold, messages must be kept in a
replica's log until it knows that the requests they concern have been
executed by at least \(f + 1\) non-faulty replicas and it can prove this
to others in view changes. In addition, if some replica misses messages
that were discarded by all non-faulty replicas, it will need to be
brought up to date by transferring all or a portion of the service
state. Therefore, replicas also need some proof that the state is
correct.

Generating these proofs after executing every operation would be
expensive. Instead, they are generated periodically, when a request with
a sequence number divisible by some constant (e.g., 100) is executed. We
will refer to the states produced by the execution of these requests as
checkpoints and we will say that a checkpoint with a proof is a stable
checkpoint.

A replica maintains several logical copies of the service state: the
last stable checkpoint, zero or more checkpoints that are not stable,
and a current state. Copy-on-write techniques can be used to reduce the
space overhead to store the extra copies of the state, as discussed in
Section 6.3.

The proof of correctness for a checkpoint is generated as follows. When
a replica i produces a checkpoint, it multicasts a message
\(\langle\text{CHECKPOINT},n,d,i\rangle_{\sigma_{i}}\) to the other
replicas, where n is the sequence number of the last request whose
execution is reflected in the state and d is the digest of the state.
Each replica collects checkpoint messages in its log until it has
\(2f + 1\) of them for sequence number n with the same digest d signed
by different replicas (including possibly its own such message). These
\(2f + 1\) messages are the proof of correctness for the checkpoint.

A checkpoint with a proof becomes stable and the replica discards all
pre-prepare, prepare, and commit messages with sequence number less than
or equal to n from its log; it also discards all earlier checkpoints and
checkpoint messages.

Computing the proofs is efficient because the digest can be computed
using incremental cryptography {[}1{]} as discussed in Section 6.3, and
proofs are generated rarely.

The checkpoint protocol is used to advance the low and high water marks
(which limit what messages will be accepted). The low-water mark h is
equal to the sequence number of the last stable checkpoint. The high
water mark \(H = h + k\) , where k is big enough so that replicas do not
stall waiting for a checkpoint to become stable. For example, if
checkpoints are taken every 100 requests, k might be 200.

\section{4.4 View Changes}\label{view-changes}

The view-change protocol provides liveness by allowing the system to
make progress when the primary fails. View changes are triggered by
timeouts that prevent backups from waiting indefinitely for requests to
execute. A backup is waiting for a request if it received a valid
request

and has not executed it. A backup starts a timer when it receives a
request and the timer is not already running. It stops the timer when it
is no longer waiting to execute the request, but restarts it if at that
point it is waiting to execute some other request.

If the timer of backup i expires in view v, the backup starts a view
change to move the system to view \(v + 1\) . It stops accepting
messages (other than checkpoint, view-change, and new-view messages) and
multicasts a
\(\langle VIEW-CHANGE, v + 1, n, C, P, i \rangle_{\sigma_i}\) message to
all replicas. Here n is the sequence number of the last stable
checkpoint s known to i, C is a set of \(2f + 1\) valid checkpoint
messages proving the correctness of s, and P is a set containing a set
\(P_m\) for each request m that prepared at i with a sequence number
higher than n.~Each set \(P_m\) contains a valid pre-prepare message
(without the corresponding client message) and 2f matching, valid
prepare messages signed by different backups with the same view,
sequence number, and the digest of m.

When the primary p of view \(v + 1\) receives 2f valid view-change
messages for view \(v + 1\) from other replicas, it multicasts a
\(\langle\) NEW-VIEW, \(v + 1\) , V, \(O\rangle_{\sigma_{p}}\) message
to all other replicas, where V is a set containing the valid view-change
messages received by the primary plus the view-change message for
\(v + 1\) the primary sent (or would have sent), and O is a set of
pre-prepare messages (without the piggybacked request). O is computed as
follows:

\begin{enumerate}
\def\labelenumi{\arabic{enumi}.}
\item
  The primary determines the sequence number min-s of the latest stable
  checkpoint in V and the highest sequence number max-s in a prepare
  message in V.
\item
  The primary creates a new pre-prepare message for view \(v + 1\) for
  each sequence number \(n\) between min- \(s\) and max- \(s\) . There
  are two cases: (1) there is at least one set in the \(\mathcal{P}\)
  component of some view-change message in \(\mathcal{V}\) with sequence
  number \(n\) , or (2) there is no such set. In the first case, the
  primary creates a new message
  \(\langle \text{PRE-PREPARE}, v + 1, n, d \rangle_{\sigma_p}\) , where
  \(d\) is the request digest in the pre-prepare message for sequence
  number \(n\) with the highest view number in \(\mathcal{V}\) . In the
  second case, it creates a new pre-prepare message
  \(\langle \text{PRE-PREPARE}, v + 1, n, d^{null} \rangle_{\sigma_p}\)
  , where \(d^{null}\) is the digest of a special null request; a null
  request goes through the protocol like other requests, but its
  execution is a no-op. (Paxos {[}18{]} used a similar technique to fill
  in gaps.)
\end{enumerate}

Next the primary appends the messages in O to its log. If min-s is
greater than the sequence number of its latest stable checkpoint, the
primary also inserts the proof of stability for the checkpoint with
sequence number min-s in its log, and discards information from the log
as discussed in Section 4.3. Then it enters view \(v+1\) : at this point
it is able to accept messages for view \(v+1\) .

A backup accepts a new-view message for view \(v + 1\) if it is signed
properly, if the view-change messages it contains are valid for view
\(v + 1\) , and if the set \(\mathcal{O}\) is correct; it verifies the
correctness of \(\mathcal{O}\) by performing a computation similar to
the one used by the primary to create \(\mathcal{O}\) . Then it adds the
new information to its log as described for the primary, multicasts a
prepare for each message in \(\mathcal{O}\) to all the other replicas,
adds these prepares to its log, and enters view \(v + 1\) .

Thereafter, the protocol proceeds as described in Section 4.2. Replicas
redo the protocol for messages between min-s and max-s but they avoid
re-executing client requests (by using their stored information about
the last reply sent to each client).

A replica may be missing some request message m or a stable checkpoint
(since these are not sent in new-view messages.) It can obtain missing
information from another replica. For example, replica i can obtain a
missing checkpoint state s from one of the replicas whose checkpoint
messages certified its correctness in V. Since \(f + 1\) of those
replicas are correct, replica i will always obtain s or a later
certified stable checkpoint. We can avoid sending the entire checkpoint
by partitioning the state and stamping each partition with the sequence
number of the last request that modified it. To bring a replica up to
date, it is only necessary to send it the partitions where it is out of
date, rather than the whole checkpoint.

\section{4.5 Correctness}\label{correctness}

This section sketches the proof that the algorithm provides safety and
liveness; details can be found in {[}4{]}.

\section{4.5.1 Safety}\label{safety}

As discussed earlier, the algorithm provides safety if all non-faulty
replicas agree on the sequence numbers of requests that commit locally.

In Section 4.2, we showed that if \(\text{prepared}(m, v, n, i)\) is
true, \(\text{prepared}(m', v, n, j)\) is false for any non-faulty
replica j (including i = j) and any \(m'\) such that \(D(m') \neq D(m)\)
. This implies that two non-faulty replicas agree on the sequence number
of requests that commit locally in the same view at the two replicas.

The view-change protocol ensures that non-faulty replicas also agree on
the sequence number of requests that commit locally in different views
at different replicas. A request m commits locally at a non-faulty
replica with sequence number n in view v only if committed(m, v, n) is
true. This means that there is a set \(R_{1}\) containing at least
\(f + 1\) non-faulty replicas such that prepared(m, v, n, i) is true for
every replica i in the set.

Non-faulty replicas will not accept a pre-prepare for view \(v' > v\)
without having received a new-view message for \(v'\) (since only at
that point do they enter the view). But any correct new-view message for
view \(v' > v\) contains correct view-change messages from every replica
i in a

set \(R_{2}\) of \(2f + 1\) replicas. Since there are \(3f + 1\)
replicas, \(R_{1}\) and \(R_{2}\) must intersect in at least one replica
k that is not faulty. k's view-change message will ensure that the fact
that m prepared in a previous view is propagated to subsequent views,
unless the new-view message contains a view-change message with a stable
checkpoint with a sequence number higher than n.~In the first case, the
algorithm redoes the three phases of the atomic multicast protocol for m
with the same sequence number n and the new view number. This is
important because it prevents any different request that was assigned
the sequence number n in a previous view from ever committing. In the
second case no replica in the new view will accept any message with
sequence number lower than n.~In either case, the replicas will agree on
the request that commits locally with sequence number n.

\section{4.5.2 Liveness}\label{liveness}

To provide liveness, replicas must move to a new view if they are unable
to execute a request. But it is important to maximize the period of time
when at least \(2f + 1\) non-faulty replicas are in the same view, and
to ensure that this period of time increases exponentially until some
requested operation executes. We achieve these goals by three means.

First, to avoid starting a view change too soon, a replica that
multicasts a view-change message for view \(v + 1\) waits for \(2f + 1\)
view-change messages for view \(v + 1\) and then starts its timer to
expire after some time T. If the timer expires before it receives a
valid new-view message for \(v + 1\) or before it executes a request in
the new view that it had not executed previously, it starts the view
change for view \(v + 2\) but this time it will wait 2T before starting
a view change for view \(v + 3\) .

Second, if a replica receives a set of \(f + 1\) valid view-change
messages from other replicas for views greater than its current view, it
sends a view-change message for the smallest view in the set, even if
its timer has not expired; this prevents it from starting the next view
change too late.

Third, faulty replicas are unable to impede progress by forcing frequent
view changes. A faulty replica cannot cause a view change by sending a
view-change message, because a view change will happen only if at least
\(f + 1\) replicas send view-change messages, but it can cause a view
change when it is the primary (by not sending messages or sending bad
messages). However, because the primary of view v is the replica p such
that \(p = v \mod |\mathcal{R}|\) , the primary cannot be faulty for
more than f consecutive views.

These three techniques guarantee liveness unless message delays grow
faster than the timeout period indefinitely, which is unlikely in a real
system.

\section{4.6 Non-Determinism}\label{non-determinism}

State machine replicas must be deterministic but many services involve
some form of non-determinism. For example, the time-last-modified in NFS
is set by reading the server's local clock; if this were done
independently at each replica, the states of non-faulty replicas would
diverge. Therefore, some mechanism to ensure that all replicas select
the same value is needed. In general, the client cannot select the value
because it does not have enough information; for example, it does not
know how its request will be ordered relative to concurrent requests by
other clients. Instead, the primary needs to select the value either
independently or based on values provided by the backups.

If the primary selects the non-deterministic value independently, it
concatenates the value with the associated request and executes the
three phase protocol to ensure that non-faulty replicas agree on a
sequence number for the request and value. This prevents a faulty
primary from causing replica state to diverge by sending different
values to different replicas. However, a faulty primary might send the
same, incorrect, value to all replicas. Therefore, replicas must be able
to decide deterministically whether the value is correct (and what to do
if it is not) based only on the service state.

This protocol is adequate for most services (including NFS) but
occasionally replicas must participate in selecting the value to satisfy
a service's specification. This can be accomplished by adding an extra
phase to the protocol: the primary obtains authenticated values proposed
by the backups, concatenates \(2f + 1\) of them with the associated
request, and starts the three phase protocol for the concatenated
message. Replicas choose the value by a deterministic computation on the
\(2f + 1\) values and their state, e.g., taking the median. The extra
phase can be optimized away in the common case. For example, if replicas
need a value that is ``close enough'' to that of their local clock, the
extra phase can be avoided when their clocks are synchronized within
some delta.

\section{5 Optimizations}\label{optimizations}

This section describes some optimizations that improve the performance
of the algorithm during normal-case operation. All the optimizations
preserve the liveness and safety properties.

\section{5.1 Reducing Communication}\label{reducing-communication}

We use three optimizations to reduce the cost of communication. The
first avoids sending most large replies. A client request designates a
replica to send the result; all other replicas send replies containing
just the digest of the result. The digests allow the client to check the
correctness of the result while reducing network

bandwidth consumption and CPU overhead significantly for large replies.
If the client does not receive a correct result from the designated
replica, it retransmits the request as usual, requesting all replicas to
send full replies.

The second optimization reduces the number of message delays for an
operation invocation from 5 to 4. Replicas execute a request tentatively
as soon as the prepared predicate holds for the request, their state
reflects the execution of all requests with lower sequence number, and
these requests are all known to have committed. After executing the
request, the replicas send tentative replies to the client. The client
waits for \(2f + 1\) matching tentative replies. If it receives this
many, the request is guaranteed to commit eventually. Otherwise, the
client retransmits the request and waits for \(f + 1\) non-tentative
replies.

A request that has executed tentatively may abort if there is a view
change and it is replaced by a null request. In this case the replica
reverts its state to the last stable checkpoint in the new-view message
or to its last checkpointed state (depending on which one has the higher
sequence number).

The third optimization improves the performance of read-only operations
that do not modify the service state. A client multicasts a read-only
request to all replicas. Replicas execute the request immediately in
their tentative state after checking that the request is properly
authenticated, that the client has access, and that the request is in
fact read-only. They send the reply only after all requests reflected in
the tentative state have committed; this is necessary to prevent the
client from observing uncommitted state. The client waits for \(2f+1\)
replies from different replicas with the same result. The client may be
unable to collect \(2f+1\) such replies if there are concurrent writes
to data that affect the result; in this case, it retransmits the request
as a regular read-write request after its retransmission timer expires.

\section{5.2 Cryptography}\label{cryptography}

In Section 4, we described an algorithm that uses digital signatures to
authenticate all messages. However, we actually use digital signatures
only for view-change and new-view messages, which are sent rarely, and
authenticate all other messages using message authentication codes
(MACs). This eliminates the main performance bottleneck in previous
systems {[}29, 22{]}.

However, MACs have a fundamental limitation relative to digital
signatures --- the inability to prove that a message is authentic to a
third party. The algorithm in Section 4 and previous
Byzantine-fault-tolerant algorithms {[}31, 16{]} for state machine
replication rely on the extra power of digital signatures. We modified
our algorithm to circumvent the problem by taking advantage of specific
invariants, e.g., the invariant that no two different requests prepare
with the same view and sequence number at two non-faulty replicas. The
modified algorithm is described in {[}5{]}. Here we sketch the main
implications of using MACs.

MACs can be computed three orders of magnitude faster than digital
signatures. For example, a 200MHz Pentium Pro takes 43ms to generate a
1024-bit modulus RSA signature of an MD5 digest and 0.6ms to verify the
signature {[}37{]}, whereas it takes only 10.3μs to compute the MAC of a
64-byte message on the same hardware in our implementation. There are
other public-key cryptosystems that generate signatures faster, e.g.,
elliptic curve public-key cryptosystems, but signature verification is
slower {[}37{]} and in our algorithm each signature is verified many
times.

Each node (including active clients) shares a 16-byte secret session key
with each replica. We compute message authentication codes by applying
MD5 to the concatenation of the message with the secret key. Rather than
using the 16 bytes of the final MD5 digest, we use only the 10 least
significant bytes. This truncation has the obvious advantage of reducing
the size of MACs and it also improves their resilience to certain
attacks {[}27{]}. This is a variant of the secret suffix method
{[}36{]}, which is secure as long as MD5 is collision resistant {[}27,
8{]}.

The digital signature in a reply message is replaced by a single MAC,
which is sufficient because these messages have a single intended
recipient. The signatures in all other messages (including client
requests but excluding view changes) are replaced by vectors of MACs
that we call authenticators. An authenticator has an entry for every
replica other than the sender; each entry is the MAC computed with the
key shared by the sender and the replica corresponding to the entry.

The time to verify an authenticator is constant but the time to generate
one grows linearly with the number of replicas. This is not a problem
because we do not expect to have a large number of replicas and there is
a huge performance gap between MAC and digital signature computation.
Furthermore, we compute authenticators efficiently; MD5 is applied to
the message once and the resulting context is used to compute each
vector entry by applying MD5 to the corresponding session key. For
example, in a system with 37 replicas (i.e., a system that can tolerate
12 simultaneous faults) an authenticator can still be computed much more
than two orders of magnitude faster than a 1024-bit modulus RSA
signature.

The size of authenticators grows linearly with the number of replicas
but it grows slowly: it is equal to
\(30 \times \left\lfloor \frac{n-1}{3} \right\rfloor\) bytes. An
authenticator is smaller than an RSA signature with a 1024-bit modulus
for \(n \leq 13\) (i.e., systems that can tolerate up to 4 simultaneous
faults), which we expect to be true in most configurations.

\section{6 Implementation}\label{implementation}

This section describes our implementation. First we discuss the
replication library, which can be used as a basis for any replicated
service. In Section 6.2 we describe how we implemented a replicated NFS
on top of the replication library. Then we describe how we maintain
checkpoints and compute checkpoint digests efficiently.

\section{6.1 The Replication Library}\label{the-replication-library}

The client interface to the replication library consists of a single
procedure, invoke, with one argument, an input buffer containing a
request to invoke a state machine operation. The invoke procedure uses
our protocol to execute the requested operation at the replicas and
select the correct reply from among the replies of the individual
replicas. It returns a pointer to a buffer containing the operation
result.

On the server side, the replication code makes a number of upcalls to
procedures that the server part of the application must implement. There
are procedures to execute requests (execute), to maintain checkpoints of
the service state (make\_checkpoint, delete\_checkpoint), to obtain the
digest of a specified checkpoint (get\_digest), and to obtain missing
information (get\_checkpoint, set\_checkpoint). The execute procedure
receives as input a buffer containing the requested operation, executes
the operation, and places the result in an output buffer. The other
procedures are discussed further in Sections 6.3 and 6.4.

Point-to-point communication between nodes is implemented using UDP, and
multicast to the group of replicas is implemented using UDP over IP
multicast {[}7{]}. There is a single IP multicast group for each
service, which contains all the replicas. These communication protocols
are unreliable; they may duplicate or lose messages or deliver them out
of order.

The algorithm tolerates out-of-order delivery and rejects duplicates.
View changes can be used to recover from lost messages, but this is
expensive and therefore it is important to perform retransmissions.
During normal operation recovery from lost messages is driven by the
receiver: backups send negative acknowledgments to the primary when they
are out of date and the primary retransmits pre-prepare messages after a
long timeout. A reply to a negative acknowledgment may include both a
portion of a stable checkpoint and missing messages. During view
changes, replicas retransmit view-change messages until they receive a
matching new-view message or they move on to a later view.

The replication library does not implement view changes or
retransmissions at present. This does not compromise the accuracy of the
results given in Section 7 because the rest of the algorithm is
completely implemented (including the manipulation of the timers that
trigger view changes) and because we have formalized the complete
algorithm and proved its correctness {[}4{]}.

\section{6.2 BFS: A Byzantine-Fault-tolerant File
System}\label{bfs-a-byzantine-fault-tolerant-file-system}

We implemented BFS, a Byzantine-fault-tolerant NFS service, using the
replication library. Figure 2 shows the architecture of BFS. We opted
not to modify the kernel NFS client and server because we did not have
the sources for the Digital Unix kernel.

A file system exported by the fault-tolerant NFS service is mounted on
the client machine like any regular NFS file system. Application
processes run unmodified and interact with the mounted file system
through the NFS client in the kernel. We rely on user level relay
processes to mediate communication between the standard NFS client and
the replicas. A relay receives NFS protocol requests, calls the invoke
procedure of our replication library, and sends the result back to the
NFS client.

\textit{[Image omitted: images/eaa6c8282ff07697665e4c6faba1e2cd6131ad63474192266eac339735e32d19.jpg]}

flowchart

\begin{Shaded}
\begin{Highlighting}[]
\NormalTok{graph TD}
\NormalTok{    A["client"] {-}{-}\textgreater{}|kernel NFS client| B["Andrew benchmark"]}
\NormalTok{    A {-}{-}\textgreater{}|kernel NFS client| C["relay replication library"]}
\NormalTok{    C {-}{-}\textgreater{}|kernel VM| D["snfsd replication library"]}
\NormalTok{    D {-}{-}\textgreater{}|replica 0| E["snfsd"]}
\NormalTok{    D {-}{-}\textgreater{}|replica n| F["snfsd"]}
\NormalTok{    F {-}{-}\textgreater{}|kernel VM| G["snfsd"]}
\NormalTok{    style A fill:\#f9f,stroke:\#333}
\NormalTok{    style B fill:\#ccf,stroke:\#333}
\NormalTok{    style C fill:\#ccf,stroke:\#333}
\NormalTok{    style D fill:\#cfc,stroke:\#333}
\NormalTok{    style E fill:\#fcc,stroke:\#333}
\NormalTok{    style F fill:\#fcc,stroke:\#333}
\NormalTok{    style G fill:\#fcc,stroke:\#333}
\end{Highlighting}
\end{Shaded}

Figure 2: Replicated File System Architecture.

Each replica runs a user-level process with the replication library and
our NFS V2 daemon, which we will refer to as snfsd (for simple nfsd).
The replication library receives requests from the relay, interacts with
snfsd by making upcalls, and packages NFS replies into replication
protocol replies that it sends to the relay.

We implemented snfsd using a fixed-size memory-mapped file. All the file
system data structures, e.g., inodes, blocks and their free lists, are
in the mapped file. We rely on the operating system to manage the cache
of memory-mapped file pages and to write modified pages to disk
asynchronously. The current implementation uses 8KB blocks and inodes
contain the NFS status information plus 256 bytes of data, which is used
to store directory entries in directories, pointers to blocks in files,
and text in symbolic links. Directories and files may also use indirect
blocks in a way similar to Unix.

Our implementation ensures that all state machine

replicas start in the same initial state and are deterministic, which
are necessary conditions for the correctness of a service implemented
using our protocol. The primary proposes the values for
time-last-modified and time-last-accessed, and replicas select the
larger of the proposed value and one greater than the maximum of all
values selected for earlier requests. We do not require synchronous
writes to implement NFS V2 protocol semantics because BFS achieves
stability of modified data and meta-data through replication {[}20{]}.

\section{6.3 Maintaining Checkpoints}\label{maintaining-checkpoints}

This section describes how snfsd maintains checkpoints of the file
system state. Recall that each replica maintains several logical copies
of the state: the current state, some number of checkpoints that are not
yet stable, and the last stable checkpoint.

snfsd executes file system operations directly in the memory mapped file
to preserve locality, and it uses copy-on-write to reduce the space and
time overhead associated with maintaining checkpoints. snfsd maintains a
copy-on-write bit for every 512-byte block in the memory mapped file.
When the replication code invokes the make\_checkpoint upcall, snfsd
sets all the copy-on-write bits and creates a (volatile) checkpoint
record, containing the current sequence number, which it receives as an
argument to the upcall, and a list of blocks. This list contains the
copies of the blocks that were modified since the checkpoint was taken,
and therefore, it is initially empty. The record also contains the
digest of the current state; we discuss how the digest is computed in
Section 6.4.

When a block of the memory mapped file is modified while executing a
client request, snfsd checks the copy-on-write bit for the block and, if
it is set, stores the block's current contents and its identifier in the
checkpoint record for the last checkpoint. Then, it overwrites the block
with its new value and resets its copy-on-write bit. snfsd retains a
checkpoint record until told to discard it via a delete\_checkpoint
upcall, which is made by the replication code when a later checkpoint
becomes stable.

If the replication code requires a checkpoint to send to another
replica, it calls the get\_checkpoint upcall. To obtain the value for a
block, snfsd first searches for the block in the checkpoint record of
the stable checkpoint, and then searches the checkpoint records of any
later checkpoints. If the block is not in any checkpoint record, it
returns the value from the current state.

The use of the copy-on-write technique and the fact that we keep at most
2 checkpoints ensure that the space and time overheads of keeping
several logical copies of the state are low. For example, in the Andrew
benchmark experiments described in Section 7, the average checkpoint
record size is only 182 blocks with a maximum of 500.

\section{6.4 Computing Checkpoint
Digests}\label{computing-checkpoint-digests}

snfsd computes a digest of a checkpoint state as part of a
make\_checkpoint upcall. Although checkpoints are only taken
occasionally, it is important to compute the state digest incrementally
because the state may be large. snfsd uses an incremental
collision-resistant one-way hash function called AdHash {[}1{]}. This
function divides the state into fixed-size blocks and uses some other
hash function (e.g., MD5) to compute the digest of the string obtained
by concatenating the block index with the block value for each block.
The digest of the state is the sum of the digests of the blocks modulo
some large integer. In our current implementation, we use the 512-byte
blocks from the copy-on-write technique and compute their digest using
MD5.

To compute the digest for the state incrementally, \(snfsd\) maintains a
table with a hash value for each 512-byte block. This hash value is
obtained by applying MD5 to the block index concatenated with the block
value at the time of the last checkpoint. When \(make\_checkpoint\) is
called, \(snfsd\) obtains the digest \(d\) for the previous checkpoint
state (from the associated checkpoint record). It computes new hash
values for each block whose copy-on-write bit is reset by applying MD5
to the block index concatenated with the current block value. Then, it
adds the new hash value to \(d\) , subtracts the old hash value from
\(d\) , and updates the table to contain the new hash value. This
process is efficient provided the number of modified blocks is small; as
mentioned above, on average 182 blocks are modified per checkpoint for
the Andrew benchmark.

\section{7 Performance Evaluation}\label{performance-evaluation}

This section evaluates the performance of our system using two
benchmarks: a micro-benchmark and the Andrew benchmark {[}15{]}. The
micro-benchmark provides a service-independent evaluation of the
performance of the replication library; it measures the latency to
invoke a null operation, i.e., an operation that does nothing.

The Andrew benchmark is used to compare BFS with two other file systems:
one is the NFS V2 implementation in Digital Unix, and the other is
identical to BFS except without replication. The first comparison
demonstrates that our system is practical by showing that its latency is
similar to the latency of a commercial system that is used daily by many
users. The second comparison allows us to evaluate the overhead of our
algorithm accurately within an implementation of a real service.

\section{7.1 Experimental Setup}\label{experimental-setup}

The experiments measure normal-case behavior (i.e., there are no view
changes), because this is the behavior

that determines the performance of the system. All experiments ran with
one client running two relay processes, and four replicas. Four replicas
can tolerate one Byzantine fault; we expect this reliability level to
suffice for most applications. The replicas and the client ran on
identical DEC 3000/400 Alpha workstations. These workstations have a 133
MHz Alpha 21064 processor, 128 MB of memory, and run Digital Unix
version 4.0. The file system was stored by each replica on a DEC RZ26
disk. All the workstations were connected by a 10Mbit/s switched
Ethernet and had DEC LANCE Ethernet interfaces. The switch was a DEC
EtherWORKS 8T/TX. The experiments were run on an isolated network.

The interval between checkpoints was 128 requests, which causes garbage
collection to occur several times in any of the experiments. The maximum
sequence number accepted by replicas in pre-prepare messages was 256
plus the sequence number of the last stable checkpoint.

\section{7.2 Micro-Benchmark}\label{micro-benchmark}

The micro-benchmark measures the latency to invoke a null operation. It
evaluates the performance of two implementations of a simple service
with no state that implements null operations with arguments and results
of different sizes. The first implementation is replicated using our
library and the second is unreplicated and uses UDP directly. Table 1
reports the response times measured at the client for both read-only and
read-write operations. They were obtained by timing 10,000 operation
invocations in three separate runs and we report the median value of the
three runs. The maximum deviation from the median was always below 0.3\%
of the reported value. We denote each operation by a/b, where a and b
are the sizes of the operation argument and result in KBytes.

arg./res. (KB)

replicated

without replication

read-write

read-only

0/0

3.35 (309\%)

1.62 (98\%)

0.82

4/0

14.19 (207\%)

6.98 (51\%)

4.62

0/4

8.01 (72\%)

5.94 (27\%)

4.66

Table 1: Micro-benchmark results (in milliseconds); the percentage
overhead is relative to the unreplicated case.

The overhead introduced by the replication library is due to extra
computation and communication. For example, the computation overhead for
the read-write 0/0 operation is approximately 1.06ms, which includes
0.55ms spent executing cryptographic operations. The remaining 1.47ms of
overhead are due to extra communication; the replication library
introduces an extra message round-trip, it sends larger messages, and it
increases the number of messages received by each node relative to the
service without replication.

The overhead for read-only operations is significantly lower because the
optimization discussed in Section 5.1 reduces both computation and
communication overheads. For example, the computation overhead for the
read-only 0/0 operation is approximately 0.43ms, which includes 0.23ms
spent executing cryptographic operations, and the communication overhead
is only 0.37ms because the protocol to execute read-only operations uses
a single round-trip.

Table 1 shows that the relative overhead is lower for the 4/0 and 0/4
operations. This is because a significant fraction of the overhead
introduced by the replication library is independent of the size of
operation arguments and results. For example, in the read-write 0/4
operation, the large message (the reply) goes over the network only once
(as discussed in Section 5.1) and only the cryptographic overhead to
process the reply message is increased. The overhead is higher for the
read-write 4/0 operation because the large message (the request) goes
over the network twice and increases the cryptographic overhead for
processing both request and pre-prepare messages.

It is important to note that this micro-benchmark represents the worst
case overhead for our algorithm because the operations perform no work
and the unreplicated server provides very weak guarantees. Most services
will require stronger guarantees, e.g., authenticated connections, and
the overhead introduced by our algorithm relative to a server that
implements these guarantees will be lower. For example, the overhead of
the replication library relative to a version of the unreplicated
service that uses MACs for authentication is only 243\% for the
read-write 0/0 operation and 4\% for the read-only 4/0 operation.

We can estimate a rough lower bound on the performance gain afforded by
our algorithm relative to Rampart {[}30{]}. Reiter reports that Rampart
has a latency of 45ms for a multi-RPC of a null message in a 10 Mbit/s
Ethernet network of 4 SparcStation 10s {[}30{]}. The multi-RPC is
sufficient for the primary to invoke a state machine operation but for
an arbitrary client to invoke an operation it would be necessary to add
an extra message delay and an extra RSA signature and verification to
authenticate the client; this would lead to a latency of at least 65ms
(using the RSA timings reported in {[}29{]}.) Even if we divide this
latency by 1.7, the ratio of the SPECint92 ratings of the DEC 3000/400
and the SparcStation 10, our algorithm still reduces the latency to
invoke the read-write and read-only 0/0 operations by factors of more
than 10 and 20, respectively. Note that this scaling is conservative
because the network accounts for a significant fraction of Rampart's
latency {[}29{]} and Rampart's results were obtained using 300-bit
modulus RSA signatures, which are not considered secure today unless the
keys used to

generate them are refreshed very frequently.

There are no published performance numbers for SecureRing {[}16{]} but
it would be slower than Rampart because its algorithm has more message
delays and signature operations in the critical path.

\section{7.3 Andrew Benchmark}\label{andrew-benchmark}

The Andrew benchmark {[}15{]} emulates a software development workload.
It has five phases: (1) creates subdirectories recursively; (2) copies a
source tree; (3) examines the status of all the files in the tree
without examining their data; (4) examines every byte of data in all the
files; and (5) compiles and links the files.

We use the Andrew benchmark to compare BFS with two other file system
configurations: NFS-std, which is the NFS V2 implementation in Digital
Unix, and BFS-nr, which is identical to BFS but with no replication.
BFS-nr ran two simple UDP relays on the client, and on the server it ran
a thin veneer linked with a version of snfsd from which all the
checkpoint management code was removed. This configuration does not
write modified file system state to disk before replying to the client.
Therefore, it does not implement NFS V2 protocol semantics, whereas both
BFS and NFS-std do.

Out of the 18 operations in the NFS V2 protocol only getattr is
read-only because the time-last-accessed attribute of files and
directories is set by operations that would otherwise be read-only,
e.g., read and lookup. The result is that our optimization for read-only
operations can rarely be used. To show the impact of this optimization,
we also ran the Andrew benchmark on a second version of BFS that
modifies the lookup operation to be read-only. This modification
violates strict Unix file system semantics but is unlikely to have
adverse effects in practice.

For all configurations, the actual benchmark code ran at the client
workstation using the standard NFS client implementation in the Digital
Unix kernel with the same mount options. The most relevant of these
options for the benchmark are: UDP transport, 4096-byte read and write
buffers, allowing asynchronous client writes, and allowing attribute
caching.

We report the mean of 10 runs of the benchmark for each configuration.
The sample standard deviation for the total time to run the benchmark
was always below 2.6\% of the reported value but it was as high as 14\%
for the individual times of the first four phases. This high variance
was also present in the NFS-std configuration. The estimated error for
the reported mean was below 4.5\% for the individual phases and 0.8\%
for the total.

Table 2 shows the results for BFS and BFS-nr. The comparison between
BFS-strict and BFS-nr shows that the overhead of Byzantine fault
tolerance for this service is low --- BFS-strict takes only 26\% more
time to run the complete benchmark. The overhead is lower than what was
observed for the micro-benchmarks because the client spends a
significant fraction of the elapsed time computing between operations,
i.e., between receiving the reply to an operation and issuing the next
request, and operations at the server perform some computation. But the
overhead is not uniform across the benchmark phases. The main reason for
this is a variation in the amount of time the client spends computing
between operations; the first two phases have a higher relative overhead
because the client spends approximately 40\% of the total time computing
between operations, whereas it spends approximately 70\% during the last
three phases.

phase

BFS

BFS-nr

strict

r/o lookup

1

0.55 (57\%)

0.47 (34\%)

0.35

2

9.24 (82\%)

7.91 (56\%)

5.08

3

7.24 (18\%)

6.45 (6\%)

6.11

4

8.77 (18\%)

7.87 (6\%)

7.41

5

38.68 (20\%)

38.38 (19\%)

32.12

total

64.48 (26\%)

61.07 (20\%)

51.07

Table 2: Andrew benchmark: BFS vs BFS-nr. The times are in seconds.

The table shows that applying the read-only optimization to lookup
improves the performance of BFS significantly and reduces the overhead
relative to BFS-nr to 20\%. This optimization has a significant impact
in the first four phases because the time spent waiting for lookup
operations to complete in BFS-strict is at least 20\% of the elapsed
time for these phases, whereas it is less than 5\% of the elapsed time
for the last phase.

phase

BFS

NFS-std

strict

r/o lookup

1

0.55 (-69\%)

0.47 (-73\%)

1.75

2

9.24 (-2\%)

7.91 (-16\%)

9.46

3

7.24 (35\%)

6.45 (20\%)

5.36

4

8.77 (32\%)

7.87 (19\%)

6.60

5

38.68 (-2\%)

38.38 (-2\%)

39.35

total

64.48 (3\%)

61.07 (-2\%)

62.52

Table 3: Andrew benchmark: BFS vs NFS-std. The times are in seconds.

Table 3 shows the results for BFS vs NFS-std. These results show that
BFS can be used in practice --- BFS-strict takes only 3\% more time to
run the complete benchmark. Thus, one could replace the NFS V2
implementation in Digital Unix, which is used daily by many users, by
BFS without affecting the latency perceived by those users. Furthermore,
BFS with the read-only optimization for the lookup operation is actually
2\% faster than NFS-std.

The overhead of BFS relative to NFS-std is not the

same for all phases. Both versions of BFS are faster than NFS-std for
phases 1, 2, and 5 but slower for the other phases. This is because
during phases 1, 2, and 5 a large fraction (between 21\% and 40\%) of
the operations issued by the client are synchronous, i.e., operations
that require the NFS implementation to ensure stability of modified file
system state before replying to the client. NFS-std achieves stability
by writing modified state to disk whereas BFS achieves stability with
lower latency using replication (as in Harp {[}20{]}). NFS-std is faster
than BFS (and BFS-nr) in phases 3 and 4 because the client issues no
synchronous operations during these phases.

\section{8 Related Work}\label{related-work}

Most previous work on replication techniques ignored Byzantine faults or
assumed a synchronous system model (e.g., {[}17, 26, 18, 34, 6, 10{]}).
Viewstamped replication {[}26{]} and Paxos {[}18{]} use views with a
primary and backups to tolerate benign faults in an asynchronous system.
Tolerating Byzantine faults requires a much more complex protocol with
cryptographic authentication, an extra pre-prepare phase, and a
different technique to trigger view changes and select primaries.
Furthermore, our system uses view changes only to select a new primary
but never to select a different set of replicas to form the new view as
in {[}26, 18{]}.

Some agreement and consensus algorithms tolerate Byzantine faults in
asynchronous systems (e.g.,{[}2, 3, 24{]}). However, they do not provide
a complete solution for state machine replication, and furthermore, most
of them were designed to demonstrate theoretical feasibility and are too
slow to be used in practice. Our algorithm during normal-case operation
is similar to the Byzantine agreement algorithm in {[}2{]} but that
algorithm is unable to survive primary failures.

The two systems that are most closely related to our work are Rampart
{[}29, 30, 31, 22{]} and SecureRing {[}16{]}. They implement state
machine replication but are more than an order of magnitude slower than
our system and, most importantly, they rely on synchrony assumptions.

Both Rampart and SecureRing must exclude faulty replicas from the group
to make progress (e.g., to remove a faulty primary and elect a new one),
and to perform garbage collection. They rely on failure detectors to
determine which replicas are faulty. However, failure detectors cannot
be accurate in an asynchronous system {[}21{]}, i.e., they may
misclassify a replica as faulty. Since correctness requires that fewer
than 1/3 of group members be faulty, a misclassification can compromise
correctness by removing a non-faulty replica from the group. This opens
an avenue of attack: an attacker gains control over a single replica but
does not change its behavior in any detectable way; then it slows
correct replicas or the communication between them until enough are
excluded from the group.

To reduce the probability of misclassification, failure detectors can be
calibrated to delay classifying a replica as faulty. However, for the
probability to be negligible the delay must be very large, which is
undesirable. For example, if the primary has actually failed, the group
will be unable to process client requests until the delay has expired.
Our algorithm is not vulnerable to this problem because it never needs
to exclude replicas from the group.

Phalanx {[}23, 25{]} applies quorum replication techniques {[}12{]} to
achieve Byzantine fault-tolerance in asynchronous systems. This work
does not provide generic state machine replication; instead, it offers a
data repository with operations to read and write individual variables
and to acquire locks. The semantics it provides for read and write
operations are weaker than those offered by our algorithm; we can
implement arbitrary operations that access any number of variables,
whereas in Phalanx it would be necessary to acquire and release locks to
execute such operations. There are no published performance numbers for
Phalanx but we believe our algorithm is faster because it has fewer
message delays in the critical path and because of our use of MACs
rather than public key cryptography. The approach in Phalanx offers the
potential for improved scalability; each operation is processed by only
a subset of replicas. But this approach to scalability is expensive: it
requires \(n > 4f + 1\) to tolerate f faults; each replica needs a copy
of the state; and the load on each replica decreases slowly with n (it
is \(O(1/\sqrt{n})\) ).

\section{9 Conclusions}\label{conclusions}

This paper has described a new state-machine replication algorithm that
is able to tolerate Byzantine faults and can be used in practice: it is
the first to work correctly in an asynchronous system like the Internet
and it improves the performance of previous algorithms by more than an
order of magnitude.

The paper also described BFS, a Byzantine-fault-tolerant implementation
of NFS. BFS demonstrates that it is possible to use our algorithm to
implement real services with performance close to that of an
unreplicated service --- the performance of BFS is only 3\% worse than
that of the standard NFS implementation in Digital Unix. This good
performance is due to a number of important optimizations, including
replacing public-key signatures by vectors of message authentication
codes, reducing the size and number of messages, and the incremental
checkpoint-management techniques.

One reason why Byzantine-fault-tolerant algorithms will be important in
the future is that they can allow systems to continue to work correctly
even when there are software errors. Not all errors are survivable; our
approach cannot mask a software error that occurs

at all replicas. However, it can mask errors that occur independently at
different replicas, including nondeterministic software errors, which
are the most problematic and persistent errors since they are the
hardest to detect. In fact, we encountered such a software bug while
running our system, and our algorithm was able to continue running
correctly in spite of it.

There is still much work to do on improving our system. One problem of
special interest is reducing the amount of resources required to
implement our algorithm. The number of replicas can be reduced by using
f replicas as witnesses that are involved in the protocol only when some
full replica fails. We also believe that it is possible to reduce the
number of copies of the state to \(f + 1\) but the details remain to be
worked out.

\section{Acknowledgments}\label{acknowledgments}

We would like to thank Atul Adya, Chandrasekhar Boyapati, Nancy Lynch,
Sape Mullender, Andrew Myers, Liuba Shrira, and the anonymous referees
for their helpful comments on drafts of this paper.

\section{References}\label{references}

{[}1{]} M. Bellare and D. Micciancio. A New Paradigm for Collision-free
Hashing: Incrementality at Reduced Cost. In Advances in Cryptology --
Eurocrypt 97, 1997.\\
{[}2{]} G. Bracha and S. Toueg. Asynchronous Consensus and Broadcast
Protocols. Journal of the ACM, 32(4), 1995.\\
{[}3{]} R. Canneti and T. Rabin. Optimal Asynchronous Byzantine
Agreement. Technical Report \#92-15, Computer Science Department, Hebrew
University, 1992.\\
{[}4{]} M. Castro and B. Liskov. A Correctness Proof for a Practical
Byzantine-Fault-Tolerant Replication Algorithm. Technical Memo
MIT/LCS/TM-590, MIT Laboratory for Computer Science, 1999.\\
{[}5{]} M. Castro and B. Liskov. Authenticated Byzantine Fault Tolerance
Without Public-Key Cryptography. Technical Memo MIT/LCS/TM-589, MIT
Laboratory for Computer Science, 1999.\\
{[}6{]} F. Cristian, H. Aghili, H. Strong, and D. Dolev. Atomic
Broadcast: From Simple Message Diffusion to Byzantine Agreement. In
International Conference on Fault Tolerant Computing, 1985.\\
{[}7{]} S. Deering and D. Cheriton. Multicast Routing in Datagram
Internetworks and Extended LANs. ACM Transactions on Computer Systems,
8(2), 1990.\\
{[}8{]} H. Dobbertin. The Status of MD5 After a Recent Attack. RSA
Laboratories' CryptoBytes, 2(2), 1996.\\
{[}9{]} M. Fischer, N. Lynch, and M. Paterson. Impossibility of
Distributed Consensus With One Faulty Process. Journal of the ACM,
32(2), 1985.\\
{[}10{]} J. Garay and Y. Moses. Fully Polynomial Byzantine Agreement for
n \textgreater3t Processors in t+1 Rounds. SIAM Journal of Computing,
27(1), 1998.\\
{[}11{]} D. Gawlick and D. Kinkade. Varieties of Concurrency Control in
IMS/VS Fast Path. Database Engineering, 8(2), 1985.\\
{[}12{]} D. Gifford. Weighted Voting for Replicated Data. In Symposium
on Operating Systems Principles, 1979.\\
{[}13{]} M. Herlihy and J. Tygar. How to make replicated data secure.
Advances in Cryptology (LNCS 293), 1988.

{[}14{]} M. Herlihy and J. Wing. Axioms for Concurrent Objects. In ACM
Symposium on Principles of Programming Languages, 1987.\\
{[}15{]} J. Howard et al.~Scale and performance in a distributed file
system. ACM Transactions on Computer Systems, 6(1), 1988.\\
{[}16{]} K. Kihlstrom, L. Moser, and P. Melliar-Smith. The SecureRing
Protocols for Securing Group Communication. In Hawaii International
Conference on System Sciences, 1998.\\
{[}17{]} L. Lamport. Time, Clocks, and the Ordering of Events in a
Distributed System. Commun. ACM, 21(7), 1978.\\
{[}18{]} L. Lamport. The Part-Time Parliament. Technical Report 49, DEC
Systems Research Center, 1989.\\
{[}19{]} L. Lamport, R. Shostak, and M. Pease. The Byzantine Generals
Problem. ACM Transactions on Programming Languages and Systems, 4(3),
1982.\\
{[}20{]} B. Liskov et al.~Replication in the Harp File System. In ACM
Symposium on Operating System Principles, 1991.\\
{[}21{]} N. Lynch. Distributed Algorithms. Morgan Kaufmann Publishers,
1996.\\
{[}22{]} D. Malkhi and M. Reiter. A High-Throughput Secure Reliable
Multicast Protocol. In Computer Security Foundations Workshop, 1996.\\
{[}23{]} D. Malkhi and M. Reiter. Byzantine Quorum Systems. In ACM
Symposium on Theory of Computing, 1997.\\
{[}24{]} D. Malkhi and M. Reiter. Unreliable Intrusion Detection in
Distributed Computations. In Computer Security Foundations Workshop,
1997.\\
{[}25{]} D. Malkhi and M. Reiter. Secure and Scalable Replication in
Phalanx. In IEEE Symposium on Reliable Distributed Systems, 1998.\\
{[}26{]} B. Oki and B. Liskov. Viewstamped Replication: A New Primary
Copy Method to Support Highly-Available Distributed Systems. In ACM
Symposium on Principles of Distributed Computing, 1988.\\
{[}27{]} B. Preneel and P. Oorschot. MDx-MAC and Building Fast MACs from
Hash Functions. In Crypto 95, 1995.\\
{[}28{]} C. Pu, A. Black, C. Cowan, and J. Walpole. A Specialization
Toolkit to Increase the Diversity of Operating Systems. In ICMAS
Workshop on Immunity-Based Systems, 1996.\\
{[}29{]} M. Reiter. Secure Agreement Protocols. In ACM Conference on
Computer and Communication Security, 1994.\\
{[}30{]} M. Reiter. The Rampart Toolkit for Building High-Integrity
Services. Theory and Practice in Distributed Systems (LNCS 938), 1995.\\
{[}31{]} M. Reiter. A Secure Group Membership Protocol. IEEE
Transactions on Software Engineering, 22(1), 1996.\\
{[}32{]} R. Rivest. The MD5 Message-Digest Algorithm. Internet RFC-1321,
1992.\\
{[}33{]} R. Rivest, A. Shamir, and L. Adleman. A Method for Obtaining
Digital Signatures and Public-Key Cryptosystems. Communications of the
ACM, 21(2), 1978.\\
{[}34{]} F. Schneider. Implementing Fault-Tolerant Services Using The
State Machine Approach: A Tutorial. ACM Computing Surveys, 22(4),
1990.\\
{[}35{]} A. Shamir. How to share a secret. Communications of the ACM,
22(11), 1979.\\
{[}36{]} G. Tsudik. Message Authentication with One-Way Hash Functions.
ACM Computer Communications Review, 22(5), 1992.\\
{[}37{]} M. Wiener. Performance Comparison of Public-Key Cryptosystems.
RSA Laboratories' CryptoBytes, 4(1), 1998.

\end{document}
